\documentclass[11pt]{article}
\usepackage[utf8]{inputenc}
\usepackage{amsmath, amssymb, amsthm}
\usepackage[margin=1in]{geometry}
\usepackage{graphicx}
\usepackage{hyperref}
\graphicspath{{figures/}}

\begin{document}

\section*{Appendix A: Derivation of the Mating System Scalar}

To explicitly derive the mating-system scalar $\kappa$ as a function of the extra-pair paternity (EPP) rate $\alpha$, we model the realized individual gametic variances $V_{km}$ and $V_{kf}$ for a population with separate sexes.

\subsection*{A.1 Population Model and Female Output}
Consider $N$ diploid adults in a constant-size population with an equal sex ratio, arranged into $N/2$ social pairs via random pair formation. For any pair $i$, the breeding values of the female $w_f^{(i)}$ and the social male $w_m^{(i)}$ are independent draws from a distribution with a mean of $1$ and an additive genetic variance $V_A$.

Female $i$ produces a clutch $K_i$ drawn from a Poisson distribution conditional on her breeding value:
\begin{equation}
    K_i \mid w_f^{(i)} \sim \mathrm{Poisson}(2 w_f^{(i)})
\end{equation}

The female's individual gamete contribution to the next generation is $k_f^{(i)} = K_i$. Utilizing the law of total variance, her total gametic variance evaluates to:
\begin{equation}
    \text{Var}(k_f^{(i)}) = \mathbb{E}[\text{Var}(K_i \mid w_f)] + \text{Var}(\mathbb{E}[K_i \mid w_f]) = \mathbb{E}[2 w_f] + \text{Var}(2 w_f) = 2 + 4 V_A
\end{equation}

\subsection*{A.2 Male Output and EPP Allocation}
Each offspring of female $i$ is independently extra-pair with probability $\alpha$ and within-pair with probability $1-\alpha$. By standard Poisson-thinning, the within-pair offspring $K_i^{\text{wp}}$ and extra-pair offspring $K_i^{\text{ep}}$ follow independent Poisson distributions conditional on $w_f^{(i)}$.

The total reproductive output of social male $i$ is the sum of his within-pair offspring and any extra-pair offspring $E_i$ he sires on other females:
\begin{equation}
    k_m^{(i)} = K_i^{\text{wp}} + E_i
\end{equation}
The variance in male output depends strictly on how extra-pair siring $E_i$ is distributed. We evaluate two bounding models.
\newline
\textbf{Model A (Breeding-Value Independent Allocation):} Extra-pair offspring are distributed uniformly across males, independent of $w_m^{(i)}$. Thus, $\mathbb{E}[E_i] = 2\alpha$ and $\text{Var}(E_i) = 2\alpha$.
\begin{align}
    \text{Var}(k_m^{(i)}) &= \text{Var}(K_i^{\text{wp}}) + \text{Var}(E_i) \nonumber \\
    &= \mathbb{E}[2(1-\alpha) w_f] + \text{Var}(2(1-\alpha) w_f) + 2\alpha \nonumber \\
    &= 2 + 4(1-\alpha)^2 V_A
\end{align}
\textbf{Model B (Breeding-Value Proportional Allocation):} Extra-pair siring success is proportional to male fitness, such that $E_i \mid w_m^{(i)} \sim \mathrm{Poisson}(2\alpha w_m^{(i)})$. Marginalizing over the independent breeding values yields:
\begin{align}
    \text{Var}(k_m^{(i)}) &= \mathbb{E}[\text{Var}(k_m \mid w_f, w_m)] + \text{Var}(\mathbb{E}[k_m \mid w_f, w_m]) \nonumber \\
    &= 2 + 4(1-\alpha)^2 V_A + 4\alpha^2 V_A \nonumber \\
    &= 2 + 4 V_A \left[(1-\alpha)^2 + \alpha^2\right]
\end{align}

\subsection*{A.3 Defining the Mating System Scalar ($\kappa$)}
The effective population size for separate sexes is parameterized by the sum of the sex-specific variances (Hill 1979):
\begin{equation}
    N_e \approx \frac{4N}{V_{km} + V_{kf}}
\end{equation}
We reorganize this expectation into the form $N_e = \frac{N}{1 + \kappa V_A}$ to isolate the mating-system scalar:
\begin{equation}
    \kappa = \frac{V_{km} + V_{kf} - 4}{2 V_A}
\end{equation}

Substituting the sum of the derived variances for Model A yields:
\begin{equation}
    \kappa^{(A)}(\alpha) = 1 + (1-\alpha)^2
\end{equation}
Substituting the sum of the derived variances for Model B yields:
\begin{equation}
    \kappa^{(B)}(\alpha) = 2(1 - \alpha + \alpha^2)
\end{equation}
Both derivations demonstrate that the demographic scaling of mating system structure is a quadratic, rather than linear, function of the extra-pair paternity rate.

\section*{References}

Hill (1979), \emph{A note on effective population size with overlapping generations.} Genetics.\\
Caballero (1994), \emph{Developments in the prediction of effective population size.} Heredity.\\
Nunney (1993), \emph{The influence of mating system and overlapping generations on effective population size.} Evolution.\\
Charlesworth (2009), \emph{Effective population size and patterns of molecular evolution and variation.} Nat Rev Genet.

\end{document}
